\begin{textsample}{POS Dim 2 – 2018\_c5 – Score -1.00 – 2018\_c5/enade\_2018\_tec\_gest\_rec\_hum\_2dfa151ea0.txt}  \label{ex:f2_pos_318}
CAPÍTULO 5 DISTRIBUIÇÃO DOS CONCEITOS 5.1 PANORAMA NACIONAL DA DISTRIBUIÇÃO DOS CONCEITOS O Conceito Enade é um indicador de qualidade que avalia os cursos por intermédio dos desempenhos dos estudantes a partir dos resultados do Enade.
Os valores possíveis vão de 1 ( pior situação ) a 5 ( melhor situação ).
Os cursos com apenas um ou sem nenhum concluinte participante não obtêm o Conceito Enade, ficando Sem Conceito ( SC ).
Na Tabela 5.1, são apresentadas a quantidade e a distribuição de cursos de Tecnologia em Gestão de Recursos Humanos participantes do Enade/2018, por faixa de conceito e Grande Região.
Enfatiza -se, mais uma vez, que a diferença entre os cursos tabulados neste capítulo e no capítulo 2 corresponde aos cursos Sem Conceito, ou seja àqueles, em princípio, sem alunos concluintes que houvessem participado da prova ou àqueles que tivessem tão somente um aluno concluinte¹⁹.
Observando -se os dados da Tabela 5.1, nota -se que, dos 599 cursos participantes, 272 ( 45,4 \% ) classificaram -se com conceito 3, o valor modal.
Esse foi também o conceito modal em todas as regiões : Norte ( 68,4 \% ), Nordeste ( 48,9 \% ), Sudeste ( 43,6 \% ), Sul ( 48,6 \% ) e Centro-Oeste ( 34,0 \% ).
O conceito 2 foi o segundo mais frequente em nível nacional ( 22,4 \%, correspondendo a 134 cursos ), e o conceito 4, o terceiro ( 20,2 \%, correspondendo a 121 cursos ).
Houve, ainda, 26 cursos ( 4,3 \% ) que receberam conceito 5, e 17 cursos ( 2,8 \% ) que receberam conceito 1.
Dos 599 cursos de Tecnologia em Gestão de Recursos Humanos, 29 ( 4,8 \% ) ficaram Sem Conceito ( SC ). ¹⁹ Esses cursos com somente um concluinte não foram considerados no capítulo 2.
A região Norte participou com 19 cursos ou 3,2 \% do total nacional.
Desses, 13 cursos receberam o conceito 3, o que equivale a 68,4 \% do total regional, além de ser o conceito modal na região, como já comentado.
Um curso ( 5,3 \% ) recebeu conceito 1, a três cursos ( 15,8 \% ) atribuiu -se conceito 2, e a dois cursos ( 10,5 \% ), conceito 4.
Nenhum curso ficou Sem Conceito ( SC ) ou recebeu conceito 5 nessa região.
A região Nordeste participou com 94 cursos ou 15,7 \% do total nacional.
Nessa região, 46 cursos ( 48,9 \% em termos regionais ) obtiveram conceito 3, o conceito modal para a região, como já comentado.
Os conceitos 2 e 4 foram atribuídos, respectivamente, a 25 e 14 cursos ( 26,6 \% e 14,9 \%, respectivamente ).
Ainda nessa região, quatro cursos ( 4,3 \% ) foram avaliados com conceito 5.
Três cursos ficaram Sem Conceito ( SC ) e outros dois receberam conceito 1.
Dos 330 cursos participantes da região Sudeste ( 55,1 \% do total nacional ), o conceito 3 foi o conceito modal, como já anteriormente assinalado, tendo sido obtido por 144 cursos ( 43,6 \% ).
O conceito 1 foi atribuído a seis cursos ( 1,8 \% ), e o conceito 2, a 74 cursos ( 22,4 \% ).
Receberam os conceitos 4 e 5, respectivamente, 71 cursos ( 21,5 \% ) e 17 cursos ( 5,2 \% ).
Nessa região, 18 cursos ficaram Sem Conceito ( SC ).
Os 109 cursos da região Sul corresponderam a 18,2 \% do total nacional.
Como já indicado, houve predominância do conceito 3, atribuído a 48,6 \% dos cursos da região ( conceito modal ), ou seja, atribuído a 53 dos 109 cursos participantes na região Sul.
Três cursos receberam conceito 1 ( 2,8 \% ), outros 17 cursos ( 15,6 \% ), o conceito 2.
Os conceitos 4 e 5 foram atribuídos a, respectivamente, 28 cursos ( 25,7 \% ) e três cursos ( 2,8 \% ).
Cinco cursos ficaram Sem Conceito ( SC ).
Já dos 47 cursos participantes na região Centro-Oeste ( 7,8 \% do total nacional ), 16 ( 34,0 \% em termos regionais ) receberam conceito 3, o conceito modal, como já fora destacado.
Cinco cursos ( 10,6 \% ) foram avaliados com conceito 1, 15 cursos ( 31,9 \% ) obtiveram conceito 2, seis cursos ( 12,8 \% ), conceito 4 e dois cursos ( 4,3 \% ), conceito 5.
Nessa região, três cursos ficaram Sem Conceito ( SC ).
No Gráfico 5.1, é apresentada a distribuição cumulativa do Conceito Enade segundo Grande Região.
Em um gráfico de distribuição cumulativa, em cada valor de abscissa ( eixo x ), o valor apresentado no eixo das ordenadas ( eixo y ) é o valor acumulado da variável.
Por exemplo, a região Nordeste ( linha verde ) e região Centro-Oeste ( linha laranja ) apresentam um pouco mais de 80,0 \% dos cursos com conceito 3 ou abaixo deste valor.
Na região Norte ( linha azul ), todos os cursos têm conceito 4 ou abaixo disso, já que a poligonal cumulativa alcança 100,0 \% para o conceito 4.
Em linhas gerais, em um gráfico de distribuição cumulativa, poligonais mais à esquerda representam Grandes Regiões com pior distribuição, e poligonais mais à direita, Grandes Regiões com melhores conceitos.
As regiões Sul ( linha rosa ) e Sudeste ( linha vermelha ) apresentam as distribuições com os melhores conceitos, ao passo que as regiões Norte ( linha azul ), Centro-Oeste ( linha laranja ) e Nordeste ( linha verde ) apresentam distribuições com menores valores. 5.2 CONCEITOS POR CATEGORIA ADMINISTRATIVA, POR MODALIDADE DE ENSINO E GRANDE REGIÃO Na Tabela 5.2, é apresentada a distribuição dos cursos de Tecnologia em Gestão de Recursos Humanos participantes do Enade/2018, por Categoria Administrativa e Modalidade de Ensino, de acordo com os conceitos por eles alcançados, segundo as Grandes Regiões brasileiras.
Dos 599 cursos participantes, dez ( 1,7 \% ) eram ministrados em Instituições Públicas, e 589 ( 98,3 \% ), em Instituições Privadas.
Quanto à Modalidade de Ensino, 70 cursos eram a Distância, e 529, Presenciais.
De acordo com as informações da Tabela 5.2, em termos nacionais, entre Instituições Públicas e Privadas, observa -se que, dos 26 cursos avaliados com conceito 5, três eram oferecidos em IES Públicas, e 23, em IES Privadas.
Dos dez cursos participantes de IES Públicas, o conceito 5 foi o valor modal, atribuído a três cursos ( 30,0 \% ).
Dos demais cursos de Instituições Públicas participantes, os conceitos 1 e 2 receberam um curso, cada ( 10,0 \%, cada ).
Dois cursos foram avaliados com conceito 3, e outros dois, com conceito 4 ( 20,0 \%, cada ).
Um curso ficou Sem Conceito ( SC ) nesta categoria. cada ).
Dois cursos foram avaliados com conceito 3, e outros dois, com conceito 4 ( 20,0 \%, cada ).
Um curso ficou Sem Conceito ( SC ) nesta categoria.
Na rede privada, o conceito modal foi 3, atribuído a 270 cursos, o correspondente a 45,8 \% dos 589 cursos da categoria.
Dos demais cursos participantes, 16 cursos ( 2,7 \% ) receberam conceito 1 e 133 cursos ( 22,6 \% ), conceito 2.
O conceito 4 foi atribuído a 119 cursos ( 20,2 \% ), e o conceito 5, a 23 cursos ( 3,9 \% ).
Nesta Categoria Administrativa, 28 cursos ( 4,8 \% ) ficaram Sem Conceito ( SC ).
Dos 529 cursos participantes da Modalidade Presencial, o conceito 3 foi o valor modal, atribuído a 248 cursos ( 45,8 \% ).
Dos demais cursos nesta Modalidade de Ensino, 13 cursos ( 2,7 \% ) receberam conceito 1, 111 cursos ( 22,6 \% ) receberam conceito 2, 108 cursos ( 20,2 \% ), conceito 4, e 21 cursos ( 3,9 \% ), conceito 5.
Nesta Modalidade de Ensino, 28 cursos ( 4,8 \% ) ficaram Sem Conceito ( SC ).
Dos 70 cursos a Distância informados, atribuíram -se quatro cursos ao conceito 1, 23 cursos ao conceito 2, 24 cursos ao conceito 3 ( conceito modal ), 13 cursos ao conceito 4 e cinco cursos ao conceito 5.
Nessa Modalidade de Ensino, um curso ficou Sem Conceito ( SC ).
Na análise por região, observa -se que, na região Norte, as Instituições Privadas participaram com 19 cursos ( 100,0 \% do total regional ), dos quais nenhum curso ficou Sem Conceito ( SC ), ou recebeu conceito 5.
O conceito modal para as Instituições Privadas na região foi 3, atribuído a 13 cursos, correspondendo a 68,4 \% dessa categoria na região.
Três cursos foram avaliados com conceito 2, e outros dois cursos, com conceito 4.
Um curso foi atribuído ao conceito 1.
As Instituições Públicas não participaram nessa combinação de categoria e região.
Dos 19 cursos oferecidos na região Norte, 18 cursos ( 94,7 \% ) eram na Modalidade Presencial e um curso ( 5,3 \% ) na Modalidade a Distância.
Dos 18 cursos participantes da Educação Presencial, o conceito 3 foi o valor modal, atribuído a 13 cursos ( 72,2 \% ).
Dos demais cursos nesta Modalidade de Ensino, um curso ( 5,6 \% ) recebeu conceito 1, dois cursos ( 11,1 \% ) receberam conceito 2, outros dois cursos ( 11,1 \% ) receberam conceito 4 e nenhum curso foi atribuído ao conceito 5.
Nessa Modalidade de Ensino, nenhum curso ficou Sem Conceito ( SC ).
O único curso informado na Modalidade a Distância, na região Norte, recebeu conceito 2.
Na região Nordeste, a rede privada concentrou 92 dos 94 cursos participantes, o equivalente a 97,9 \% do total da região.
Entre os cursos oferecidos por IES Privadas no Nordeste, 45 cursos foram avaliados com conceito 3, o conceito modal, 25 cursos obtiveram conceito 2, 13 cursos receberam conceito 4 e quatro cursos ficaram com conceito 5.
Nessa combinação de categoria e região, três cursos ficaram Sem Conceito ( SC ) e dois cursos receberam conceito 1.
As Instituições Públicas da região Nordeste participaram com dois cursos ( 2,1 \% do total da região ).
Desses, um obteve conceito 3 e outro conceito 4.
Nenhum curso ficou Sem Conceito ( SC ) ou recebeu os demais conceitos nessa combinação de categoria e região.
Dos 88 cursos participantes da Educação Presencial, o conceito 3 foi o valor modal, atribuído a 43 cursos ( 48,9 \% ).
Dos demais cursos nessa Modalidade de Ensino, dois cursos ( 2,3 \% ) receberam conceito 1, 23 cursos ( 26,1 \% ) receberam conceito 2, 13 cursos ( 14,8 \% ) receberam conceito 4 e quatro cursos foram avaliados com conceito 5.
Nessa Modalidade de Ensino, três cursos ficaram Sem Conceito ( SC ).
Dos seis cursos informados na Modalidade a Distância, dois cursos ( 33,3 \% ) receberam conceito 2, três cursos ( 50,0 \% ) receberam conceito 3, o conceito modal, e um curso ( 16,7 \% ) recebeu conceito 4.
Nenhum curso ficou Sem Conceito ( SC ) ou recebeu os demais conceitos nessa Modalidade de Ensino e região.
Na região Sudeste, a proporção de cursos da rede privada foi de 98,2 \%, correspondendo a 324 dos 330 cursos participantes.
Nessa categoria e região, o conceito modal foi 3, atribuído a 143 cursos ( 44,1 \% ).
Ficaram Sem Conceito ( SC ), 18 cursos.
Os demais cursos foram avaliados com conceito 1 ( seis cursos ), conceito 2 ( 73 cursos ), conceito 4 ( 70 cursos ) e conceito 5 ( 14 cursos ).
Dos seis cursos oferecidos por Instituições Públicas na região Sudeste, o conceito modal foi 5, atribuído a três cursos.
Nessa região, os conceitos 2, 3 e 4 receberam um curso, cada.
Nenhum curso ficou Sem Conceito ( SC ) ou recebeu conceito 1.
Na região Sudeste, dos 330 cursos participantes, 293 foram oferecidos na Modalidade Presencial.
Para esses, a moda foi o conceito 3, com 132 cursos.
Os demais cursos presenciais da região foram alocados aos conceitos 1 ( quatro cursos ), 2 ( 60 cursos ), 4 ( 66 cursos ) e 5 ( 14 cursos ).
Ficaram Sem Conceito ( SC ) 17 cursos.
Já os 37 cursos oferecidos na Modalidade Educação a Distância obtiveram os conceitos 1 ( dois cursos ), 2 ( 14 cursos, conceito modal ), 3 ( 12 cursos ), 4 ( cinco cursos ) e 5 ( três cursos ).
Um curso ficou Sem Conceito ( SC ) na Modalidade Educação a Distância na região Sudeste.
As Instituições Privadas concentraram 108 dos 109 cursos participantes da região Sul, 99,1 \% do total regional.
Desses, 53 cursos obtiveram conceito 3, o conceito modal.
Nessa combinação de Categoria Administrativa e Grande Região, dois cursos receberam conceito 1, 17 cursos receberam conceito 2, outros 28, conceito 4.
Três cursos foram avaliados com conceito 5 e cinco cursos ficaram Sem Conceito ( SC ).
As Instituições Públicas na região Sul participaram com um curso ( 0,9 \% dos cursos da região ), avaliado com conceito 1.
Nenhum curso da região Sul oferecido por IES Públicas ficou Sem Conceito ( SC ) ou recebeu os demais conceitos.
Dos 86 cursos participantes da Educação Presencial na região Sul, o conceito 3 foi o valor modal, atribuído a 45 cursos ( 52,3 \% ).
Dos demais cursos nessa Modalidade de Ensino, dois cursos ( 2,3 \% ) receberam conceito 1, 11 cursos ( 12,8 \% ) receberam conceito 2, 22 cursos ( 25,6 \% ) receberam conceito 4 e um curso foi avaliado com conceito 5.
Nessa Modalidade de Ensino, cinco cursos ficaram Sem Conceito ( SC ).
Dos 23 cursos informados na Modalidade a Distância, a um curso ( 4,3 \% ) foi atribuído o conceito 1, seis cursos ( 26,1 \% ) receberam conceito 2, oito cursos ( 34,8 \% ) receberam conceito 3, seis cursos ( 26,1 \% ) foram avaliados com conceito 4 e dois cursos ( 8,7 \% ) com conceito 5.
Nenhum curso ficou Sem Conceito ( SC ).
Na região Centro-Oeste, 47 cursos participantes eram de Instituições Privadas ( 97,9 \% em termos regionais ).
Desses, 16 cursos receberam conceito 3, o conceito modal.
Dos demais cursos, cinco receberam conceito 1, 15 receberam conceito 2, seis receberam conceito 4 e dois cursos foram avaliados com conceito 5.
Dois cursos ficaram Sem Conceito ( SC ).
As Instituições Públicas, nessa região, participaram com apenas um curso, o qual ficou Sem Conceito ( SC ).
Dos 44 cursos participantes da Educação Presencial na região Sul, o conceito 2 e 3 receberam 15 cursos, cada ( 34,1 \%, cada ).
Dos demais cursos nessa Modalidade de Ensino, quatro cursos ( 9,1 \% ) receberam conceito 1, cinco cursos ( 11,4 \% ) receberam conceito 4 e dois cursos ( 4,5 \% ) foram avaliados com conceito 5.
Nessa Modalidade de Ensino, três cursos ficaram Sem Conceito ( SC ).
Dos três cursos informados na Modalidade a Distância, os conceitos 1, 3 e 4 foram atribuídos a um curso, cada ( 33,3 \%, cada ).
Os demais conceitos não foram atribuídos, e nenhum curso ficou Sem Conceito ( SC ).
No Gráfico 5.2, apresenta -se a distribuição do Conceito Enade, segundo a Categoria Administrativa da IES.
Os cursos em IES Públicas ( linha azul ) apresentam uma poligonal mais à direita do que os em IES Privadas ( linha verde ) e, consequentemente, uma distribuição de Conceitos Enade com valores maiores. 5.3 CONCEITOS POR ORGANIZAÇÃO ACADÊMICA E GRANDE REGIÃO Na Tabela 5.3, encontra -se a distribuição dos conceitos atribuídos aos cursos participantes do Enade/2018, na Área de Tecnologia em Gestão de Recursos Humanos, por Organização Acadêmica, segundo as Grandes Regiões brasileiras.
Dos 599 cursos de Tecnologia em Gestão de Recursos Humanos participantes, 145 eram oferecidos em Universidades, 159 em Centros Universitários, 291 em Faculdades e quatro em CEFET/IFET.
Essa distribuição corresponde a, respectivamente, 24,2 \%, 26,5 \%, 48,6 \% e 0,7 \% dos cursos.
De acordo com os dados apresentados, de todos os 26 cursos avaliados com conceito 5, cinco eram vinculados a Universidades.
Os cursos vinculados a esse tipo de Organização Acadêmica tiveram o conceito 3 como conceito modal, atribuído a 71 cursos ( 49,0 \% ).
Os demais cursos vinculados a Universidades receberam os conceitos 1 ( quatro cursos ), 2 ( 40 cursos ), 4 ( 20 cursos ) e conceito 5 ( cinco cursos ).
Ficaram Sem Conceito ( SC ), cinco cursos.
Dos 159 cursos oferecidos em Centros Universitários, o conceito modal também foi 3, atribuído a 71 cursos ( 44,7 \% ).
Oito cursos vinculados a esse tipo de Organização Acadêmica ficaram Sem Conceito ( SC ), e os outros cursos receberam os conceitos 1 ( quatro cursos ), 2 ( 32 cursos ), 4 ( 38 cursos ) e 5 ( seis cursos ).
Dos 291 cursos mantidos por Faculdades, 130 ( 44,7 \% ) receberam conceito 3, o conceito modal.
Seguido dos conceitos 2 e 4, que receberam 62 cursos, cada.
Na sequência, 13 cursos obtiveram conceito 5, nove cursos, conceito 1, e 15 cursos ficaram Sem Conceito ( SC ).
Dos 3 cursos mantidos por CEFET/IFET, dois cursos ( 50,0 \% ) receberam conceito 5, o conceito modal, e um curso recebeu conceito 4.
Nesse tipo de Organização Acadêmica um curso ficou Sem Conceito ( SC ) e nenhum curso recebeu os demais conceitos.
Considerando -se separadamente as regiões brasileiras, verifica -se que, na região Norte, as Universidades participaram com apenas três cursos dos 19 avaliados.
Os três cursos informados receberam conceito 3.
Nenhum curso ficou Sem Conceito ( SC ) ou recebeu os demais conceitos nessa combinação de categoria e região.
Os Centros Universitários da região Norte foram representados por quatro cursos, os conceitos 1 e 2 receberam um curso, cada, e o conceito 3, dois cursos.
Nenhum curso ficou Sem Conceito ( SC ) ou recebeu os demais conceitos nessa combinação de categoria e região.
As Faculdades participaram com 12 cursos na região Norte, oito cursos foram avaliados com conceito 3, o conceito modal, e os demais cursos receberam os conceitos 2 ( dois cursos ) e 4 ( dois cursos ).
Nenhum curso ficou Sem Conceito ( SC ) ou recebeu os conceitos 1 ou 5.
Já os CEFET/IFET não participaram com cursos na região Norte.
Na região Nordeste, as Universidades participaram com 13 dos 94 cursos da Área de Tecnologia em Gestão de Recursos Humanos oferecidos.
Apenas um desses 13 cursos ficou Sem Conceito ( SC ).
O conceito modal 3 foi atribuído seis cursos.
Os demais cursos obtiveram conceitos 2 ( quatro cursos ) e 4 ( dois cursos ).
Um curso oferecido por Universidades no Nordeste ficou sem conceito e os conceitos 1 e 5 não tiveram cursos atribuídos.
Os Centros Universitários contaram com 29 cursos participantes na região Nordeste, 17 dos quais receberam o conceito modal 3.
Para os demais cursos, foram atribuídos os conceitos 2 ( sete cursos ) e 4 ( quatro cursos ).
Um curso oferecido por Centros Universitários no Nordeste ficou Sem Conceito ( SC ) e os conceitos 1 e 5 não foram atribuídos.
As Faculdades foram representadas por 51 cursos na região Nordeste, 23 dos quais receberam conceito 3, o conceito modal.
Seguido do conceito 2, com 14 cursos.
Os demais conceitos receberam os conceitos 1 ( dois cursos ), 4 ( sete cursos ) e 5 ( quatro cursos ).
Na região Nordeste, um curso mantido por Faculdades ficou Sem Conceito ( SC ).
Os CEFET/IFET participaram com um curso na região Nordeste avaliado com conceito 4.
Nenhum curso ficou Sem Conceito ( SC ) ou recebeu os demais conceitos nesse tipo de Organização Acadêmica.
Na região Sudeste, as Universidades concentraram 102 dos 330 cursos de Tecnologia em Gestão de Recursos Humanos da região.
Dos cursos oferecidos em Universidades, nessa região, o conceito modal foi 3, atribuído a 51 cursos.
Os demais cursos receberam os conceitos 1 ( dois cursos ), 2 ( 31 cursos ), 4 ( 11 cursos ) e 5 ( quatro cursos ).
Três cursos ficaram Sem Conceito ( SC ).
Os Centros Universitários participaram com 89 cursos na região Sudeste, dos quais 35 obtiveram o conceito modal 3.
O conceito 1 foi atribuído três cursos ; o conceito 2, a 15 cursos.
O conceito 4 foi atribuído a 25 cursos, e o conceito 5, a quatro.
Sete cursos ficaram Sem Conceito ( SC ) nessa combinação de Organização Acadêmica e região.
Já as Faculdades No Gráfico 5.3, apresenta -se a distribuição do Conceito Enade segundo a Organização Acadêmica da IES.
Os cursos em CEFET/IFET ( linha laranja ) apresentam uma poligonal mais à direita do que as demais, denotando uma melhor distribuição de Conceitos Enade.
No outro extremo, encontram -se os cursos mantidos por Universidades ( linha azul ), apresentando a respectiva poligonal mais à esquerda e, consequentemente, distribuição de Conceitos Enade com menores valores, seguido dos cursos oferecidos em Centros Universitários ( linha verde ) e Faculdades ( linha vermelha ), que apresentam distribuições em situação intermediária e semelhante entre si. foram representadas por 137 cursos na região Sudeste, dos quais oito ficaram Sem Conceito ( SC ).
O conceito modal 3 foi atribuído a 58 cursos.
Aos demais cursos foram atribuídos os conceitos 1 ( um curso ), 2 ( 28 cursos ), 4 ( 35 cursos ) e 5 ( sete cursos ).
Os CEFET/IFET participaram com dois cursos avaliados com conceito 5.
Nesse tipo de Organização Acadêmica nenhum curso ficou Sem Conceito ( SC ) ou recebeu os demais conceitos Dos 109 cursos da região Sul, 24 cursos eram mantidos por Universidades, para os quais o conceito modal foi 3, atribuído a 11 cursos.
Os demais cursos receberam os conceitos 1 ( um curso ), 2 ( quatro cursos ) e 4 ( sete cursos ).
Nesse tipo de organização, um curso da região Sul ficou Sem Conceito ( SC ) ou recebeu conceito 5.
Dos 26 cursos participantes de Centros Universitários da região Sul, nenhum ficou Sem Conceito ( SC ), ou recebeu conceito 1.
Aos demais cursos foram atribuídos os conceitos 2 ( três cursos ), 3 ( 14 cursos, conceito modal ), 4 ( oito cursos ) e 5 ( um curso ).
Quanto aos 58 cursos vinculados a Faculdades na região Sul, 28 receberam o conceito modal 3.
O conceito 2 foi atribuído a dez cursos e o conceito 4, a 13 cursos.
O conceito 1 e 5 receberam dois cursos, cada.
Nesse tipo de Organização Acadêmica e região, quatro cursos ficaram Sem Conceito ( SC ).
Já os CEFET/IFET não participaram com cursos nessa região.
Na região Centro-Oeste, três dos 47 cursos oferecidos eram mantidos por Universidades.
Para este tipo de Organização Acadêmica, nenhum curso ficou Sem Conceito ( SC ) ou recebeu os conceitos 3 ou 4.
Os conceitos 1, 2 e 5 receberam um curso, cada.
Os Centros Universitários da região Centro-Oeste contaram com 11 cursos.
Seis cursos receberam o conceito modal 2.
Os demais receberam os conceitos 3 ( três cursos ), 4 ( um curso ) e 5 ( um curso ).
Nenhum curso ficou Sem Conceito ( SC ) ou recebeu conceito 1.
Dos 32 cursos oferecidos em Faculdades na região Centro-Oeste, 13 receberam o conceito modal 3, quatro cursos obtiveram conceito 1, oito cursos, conceito 2 e cinco cursos, conceito 4.
Nesta região, dois cursos ficaram Sem Conceito ( SC ) e nenhum curso obteve conceito 5.
Já os CEFET/IFET participaram apenas com um curso, o qual ficou Sem Conceito ( SC ).
No Gráfico 5.3, apresenta -se a distribuição do Conceito Enade segundo a Organização Acadêmica da IES.
Os cursos em CEFET/IFET ( linha laranja ) apresentam uma poligonal mais à direita do que as demais, denotando uma melhor distribuição de Conceitos Enade.
No outro extremo, encontram -se os cursos mantidos por Universidades ( linha azul ), apresentando a respectiva poligonal mais à esquerda e, consequentemente, distribuição de Conceitos Enade com menores valores, seguido dos cursos oferecidos em Centros Universitários ( linha verde ) e Faculdades ( linha vermelha ), que apresentam distribuições em situação intermediária e semelhante entre si.

% matched lemmas: 
\end{textsample}
